

\chapter{Introduction}


The International Civil Aviation Organization (ICAO) lists runway incursions as one of the five global high risk categories, and the deadly accident earlier this year at Tokyo Haneda Airport claimed the lives of 5 Japanese Coast Guard members highlights a gap within aviation safety. Despite Haneda Airport being equipped with a surface movement radar, it ultimately failed to detect the runway incursion. American airports are getting equipped with Airport Surface Detection Equipment (ASDE), and while ASDE has been able to reduce the number of operational incidents, pilots now need a flight deck based solution that will provide them with the relevant information to react in when required.

\section{Background}


The risk of runway incursions is only bound to increase as the volume of air traffic rises, and while aircraft have Airborne Collision Avoidance System (ACAS) and Traffic Collision Avoidance System (TCAS) to mitigate the risk of in-air collisions, runway safety measures are relatively lacking. Separation on the runway is dependent on operational procedures by the local air traffic controller (ATC), and the only indicator for runway occupancy that can be seen from the flight deck are the runway status lights (RWSL). While airports around the world are installing surface movement guidance and control (SMGC) equipment such as ground radar to track aircraft and vehicle movement throughout the airport, this only displays the information and provides alerts to the controller on the ground. 

A flight deck based system, much like ACAS and TCAS, can provide the information to the pilots in real-time so that they can act accordingly. America's Federal Aviation Authority (FAA) recently certified Garmin's Runway Occupancy Awareness (ROA) for the G1000, an avionics suite designed for general aviation. Garmin's ROA is the first certified use of Enhanced Traffic Situational Awareness on the Airport Surface with Indications and Alerts (SURF IA). Garmin expects ROA certification for the G5000 in late 2024, which is designed for transport category aircraft. Honeywell, a EUROCONTROL partner, has sucessfully demonstrated a runway surface alert system on a Boeing 757, and is hoping for FAA certification by 2026. However, as to date, there are no certified SURF IA systems that are in regular operation in commercial aviation. 

While SURF IA seems to be a promising solution to runway incursions, this is an extremely new technology, with no precedent. The industry standard set by the Radio Technical Commission for Aeronautics (RTCA), the DO-323, has no data on its effectiveness on detecting runway incursions. 


\section{Scope and Delimitations}
\label{sec:bilder}

This study will cover (500+) runway incursions worldwide dating back to 1995. This study will make use of the FAA’s Runway Incursion Severity (RIS) model and only focus on Category A and B incidents, which FAA considers to be serious runway incursions, and accidents. ICAO's Annex 11 defines Category A as a serious incident where collision was narrowly avoided, and Category B as an incident where separation decreases and there is significant potential for a collision, which may result in time critical corrective/evasive response to avoid a collision. Accidents are runway incursions that resulted in collisions.

This study will identify and characterize the aircraft types, airports, states and runway geometry at time of convergence, causes, and outcomes, and will solely focus on whether the incident will be detected by the SURF IA system or not.


%\vspace{0,5cm}
%$ $ 
%If a figure ends the section, the spacing to the next heading may be too small. The spacing can be adjusted by setting a space in the math environment ($ $) or using the command \vspace{0.5cm}.




\section{Limitations}

This study will not take into account human machine interactions, such as the pilot's reaction time to the alert, weather conditions, if the alerts and warnings displayed are sufficient for the pilots to recognize a conflict scenario, or interactions within the flight deck, such as pilot head down time. This study will solely focus on the effectiveness of SURF IA detection logic when presented with historical conflict scenarios.

This study will primarily focus on commercial aircraft, as it is assumed that only ADSB equipped aircraft can be identified as traffic by the SURF IA system, and that general aviation aircraft and ground vehicles do not generally carry ADSB transponders. While general aviation aircraft and ground vehicles are viable targets for surface movement radar, this information is not displayed on the flight deck and is out of the scope of this study. The team also assumes that the SURF IA system installed in the aircraft is in perfect operational condition, with 100 percent data integrity and accuracy with regards to time and position.

\section{Objectives}

This study's objective is to identify, characterize, and classify runway incursions. Runway layouts, conflict geometry, causes, and aircraft maneuvers are to be analyzed and frequently occurring parameter sets that lead to runway incursions are to be identified.

This study will then assess the effectiveness of the industry standard by applying its alerting behaviors to the data sets to determine the effectiveness of detecting parameters that may lead to runway incursions.

This study's ultimate objective is to provide recommendations to increase the alert effectiveness of the industry standard.

\newpage
